\section{结论与展望}
\label{sec:conclusion}

高密度城际高速铁路作为我国综合交通体系的关键骨干，其运行安全、准点与效率高度依赖于实时重调度系统的智能性与鲁棒性。然而，现有方法在应对“多车同区间”（$M > N$）这一典型高密度场景下的复杂扰动时，普遍存在模型失配、抗扰能力单一、策略割裂及实时性不足等瓶颈。针对上述挑战，本文以先进控制理论为基石，围绕状态空间建模—鲁棒控制—迭代学习—迭代模糊融合四大核心环节，构建了一套兼具理论严谨性、算法实时性与工程实用性的列车实时重调度控制新范式。主要结论如下：

第一，建立了面向高密度运行特征的精确离散状态空间模型。针对城际高铁“多车同区间”（$M > N$）这一典型高密度场景，突破传统“单区间单车”的简化建模范式，首次构建了能够完整刻画多列车在共享运行区间内动态交互关系的离散时间状态空间模型。该模型显式纳入乘客上下车引起的停站时间扰动、列车追踪间隔约束以及延误在时空维度上的传播机制，准确反映列车群在密集追踪条件下的耦合动力学特性。通过引入批处理结构，进一步支持预测时域与控制时域的协同优化，为后续先进控制器的设计奠定了兼具物理真实性与计算可实现性的系统描述基础。

第二，提出了基于 Tube-MPC 框架的鲁棒模型预测控制器（RMPC）。面向极端天气、信号故障或临时限速等突发性、有界但未知的非重复性扰动，设计了一种基于管状模型预测控制（Tube-MPC）的鲁棒调度策略。通过构造名义系统轨迹与误差动态的鲁棒正不变集，并设计辅助状态反馈律，将所有可能的实际系统轨迹严格约束在以名义轨迹为中心的安全“管”内。该方法在保证最小行车间隔、到发时刻窗口等硬性运营约束的同时，有效隔离了最坏情形扰动对调度安全的影响，显著提升了系统在强不确定性环境下的运行韧性，克服了传统确定性 MPC 因忽略扰动边界而导致的可行性丧失或性能退化问题。

第三，设计了面向周期化运营的迭代学习控制机制（ILC）。充分利用城际高速铁路日复一日高度重复的运行规律，特别是早晚高峰客流模式的半稳态周期性特征，构建了基于批次更新的迭代学习控制架构。通过建立相邻运行周期间的误差传递映射关系，控制器能够在每次运营结束后利用历史跟踪误差信息修正下一轮的控制输入序列，从而实现对周期性停站延误的渐进式补偿与消除。理论分析表明，在初始状态一致且扰动变化率有界的合理假设下，闭环系统状态误差序列可收敛至一个与扰动变化幅度相关的紧致邻域，充分展现出“运行次数越多、调度精度越高”的自主进化能力，为提升长期准点率提供了有效路径。

第四，提出了融合迭代学习机制的模糊预测控制架构——迭代学习模糊预测控制（ILFPC）。基于 Takagi–Sugeno (T-S) 模糊模型，通过定义客流强度、延误传播速度等关键指标的模糊子集（NB, NM, NS, PS, PM, PB），构建12条模糊规则，动态识别当前运行工况（如平峰、高峰或异常事件），并在线调整预测控制参数。在此基础上，将迭代学习机制嵌入模糊预测控制框架，使控制器能够利用历史运行数据逐轮优化控制输入，实现对重复性扰动的渐进消除与非重复性扰动的自适应响应。仿真结果表明，该方法在首次运行中即可将大客流扰动下的最大延误控制在9.2分钟以内，并随着迭代持续降低总延误时间，其调度性能显著优于传统的迭代学习控制、混合整数规划及鲁棒模型预测控制等基准方法。。

综上，本文所构建的 ILFPC 框架，成功实现了从被动抗扰（鲁棒控制）到主动学习（迭代学习）再到智能融合（模糊自适应）的跨越，不仅在理论上解决了高密度场景下模型适配性、扰动双重性与策略协同性三大难题，更在工程上验证了其在秒级响应、准点提升与运行平稳性方面的综合优势，为智能高铁调度控制提供了可落地的新路径。

展望未来，本研究可在以下方向进一步深化与拓展：

（1）网络化扩展：将当前单线模型推广至大规模高速铁路网，考虑枢纽换乘、跨线列车、供电分区及咽喉区资源冲突等复杂耦合约束，发展分布式/分层式 ILFPC 架构，实现路网级协同调度；

（2）感知–决策闭环：深度融合实时多源客流感知数据（如AFC刷卡、视频AI计数、手机信令），构建在线扰动辨识模块，实现模糊规则库的动态更新与控制策略的前馈补偿，提升系统对突发扰动的预判与响应能力；

（3）数据–机理混合建模：探索深度学习（如图神经网络、Transformer）与Takagi–Sugeno模糊模型的融合机制，利用数据驱动方法自动挖掘高维非线性映射关系，提升模糊规则生成的智能化水平与泛化能力；

（4）车–地–云一体化控制：在自动驾驶列车与智能调度中心协同背景下，将 ILFPC 与车载节能驾驶曲线优化、再生制动能量管理相结合，构建“调度–控制–能效”联合优化框架，推动轨道交通向绿色智能演进；

（5）安全增强机制：针对未来开放通信环境下的潜在网络攻击风险，研究具有攻击检测与弹性恢复能力的鲁棒迭代学习架构，确保控制系统在信息物理融合环境中的可信性与韧性。

上述研究方向的持续推进，有望推动列车运行控制从“经验驱动、被动响应”向“模型–数据双驱动、主动预测–自适应优化”范式的根本转变，为构建下一代高安全、高效率、高韧性、高智能的中国式现代化轨道交通系统提供核心理论支撑与关键技术引擎。
